%!TEX root = ./../../springer-mono-format/book.tex
%!TEX root = ./../../thesis-trier.tex

\Cref{prelim:thm::one-attractor-idempotent} lists several characterizations for idempotents
whose duplicate set contains at most one point,
and \cref{prelim:prop:one-attractor-idempotent-fin} 
provides an additional characterization for finite domains.
Then, \cref{prelim:trans:idempt-one-point:prop:comp}, which we need in Part~\ref{part:aut},
gives conditions when the composition of two idempotent one-point singulars
yields an idempotent one-point singular. As a consequence, we obtain
that if the composition of two simple idempotents yields a simple idempotent,
then it equals one of the two original simple idempotents.

A simple idempotent is an idempotent one-point singular.
Rystsov~\cite{Rystsov2000} introduced simple idempotents
and showed~\cite[Lem.~1]{Rystsov2000} that a simple idempotent $t \colon [n] \to [n]$
on a finite set fulfills $\applyfunc{A}{t^{-1}} \subseteq A$
or $A \subseteq \applyfunc{A}{t^{-1}}$ for every $A \subseteq [n]$.
In~\crefitem{prelim:thm::one-attractor-idempotent}{prelim:prop:item:inverse_super_or_subset2} 
we extend this result by characterizing the transformations that fulfill 
this equation as precisely the idempotents whose
duplicate set contains at most one point.

\begin{theorem}
	\label{prelim:thm::one-attractor-idempotent}
	Let $t \colon X \to X$ with~$X$ non-empty. Then the following are equivalent.
	\begin{enumerate}
	\item \label{prelim:prop:item:one-attractor} 
	$t$ is an idempotent one-point singular or equals the identity transformation.

    \item \label{prelim:prop:item:equation}
    $t = [a_{m-1}\mapsto a_m] \cdots [a_1 \mapsto a_m]$
    for $\move{t} = \{ a_1, \ldots, a_m \}$ and $\dupl{t} = \{a_m\}$ or $t = \operatorname{id}_X$.

    \item \label{prelim:prop:item:one-attractor-a} 
    $\dupl{t}$ contains at most one point and $t$ is idempotent.

	\item \label{prelim:prop:item:oneattractor-single} $\dupl{t}$ contains at most one point
	and $|\applyfunc{z}{t^{-1}}| = 1$ implies $\applyfunc{z}{t^{-1}} = \{z\}$
	for all $z \in X$.
		
	\item  \label{prelim:prop:item:invimg1}  There exists $y \in X$ such that for all $z \in X$
	we have $\applyfunc{z}{t^{-1}} = \{z\}$
	or $z \in \applyfunc{y}{t^{-1}}$.
	
	\item  \label{prelim:prop:item:invimg1a} There exists $y \in X$
	such that $x \in t^{-1}(z)$ implies $t(x) = x$ for $z,x\in X$ with $z \ne y$.	
	
	\item  \label{prelim:prop:item:mapping} For all $x,z \in X$ 
	with $t(x) \ne x$ and $t(z) \ne z$
	we have $t(x) = t(z)$, i.e., $\move{t} \subseteq t^{-1}(y)$ for some $y \in X$.
	
	\item \label{prelim:prop:item:invimg2} There exists $y \in X$ such that for all $z \in X$
	we have $\applyfunc{z}{t^{-1}} \subseteq \{z\}$
	or $z \in \applyfunc{y}{t^{-1}}$.
	
	\item \label{prelim:prop:item:excluded2} There exists $y \in X$ such that, for all $z \in t(X)$, we have
	\[
	\applyfunc{z}{t^{-1}} = \begin{cases}
		\excl{t} \cup \{y\} & \text{if } z = y, \\
		\{ z \}             & \text{otherwise.}
	\end{cases}
	\]
	
	\item \label{prelim:prop:item:inverse_super_or_subset2} For all $A \subseteq X$ we have
	$A \subseteq \applyfunc{A}{t^{-1}}$ or
	$\applyfunc{A}{t^{-1}} \subseteq A$.
\end{enumerate}
\end{theorem}
\begin{proof}
	$\ref{prelim:prop:item:one-attractor} \Rightarrow~\ref{prelim:prop:item:equation}$
	For $t \ne \operatorname{id}_X$,
	let $\dupl{t} = \{a_m\}$ and $\move{t} = \{a_1, \ldots, a_{m-1}\}$.
	Then we have $t = [a_{m-1}\mapsto a_m] \cdots [a_1 \mapsto a_m]$.
		
	$\ref{prelim:prop:item:equation} \Rightarrow~\ref{prelim:prop:item:one-attractor-a}$:
	This is clear.
	
	$\ref{prelim:prop:item:one-attractor-a} \Rightarrow~\ref{prelim:prop:item:oneattractor-single}$: 
	By assumption, $\dupl{t}$ contains at most one point. By idempotency,
	the rest of the claim follows from \crefitem{prelim:prop:idempotent}{prop:idemp:item:invy}.	
	
	$\ref{prelim:prop:item:oneattractor-single} \Rightarrow~\ref{prelim:prop:item:invimg1}$: 
	For $\dupl{t} \ne \emptyset$, let $y$
	be the unique point in $\dupl{t}$,
	otherwise, let $y \in X$ be any point.
	Next, let $z \in X$ with $t(z) \ne y$. 
	Then as $t(z) \notin \dupl{t}$ we have $t^{-1}(t(z)) = \{z\}$.
	By assumption $|t^{-1}(t(z))| = 1$
	gives $t^{-1}(t(z)) = \{t(z)\}$.
	Thus $t(z) = z$, which gives $t^{-1}(z) = \{z\}$. 

$\ref{prelim:prop:item:invimg1} \Rightarrow~\ref{prelim:prop:item:invimg1a}$: Let $x,z \in X$ with $z \ne y$ and suppose $x \in t^{-1}(z)$.
Since $t^{-1}(z) = \{z\}$, we get $x = z$ and so $t(x) = x$.

$\ref{prelim:prop:item:invimg1a} \Rightarrow~\ref{prelim:prop:item:mapping}$: Suppose $t(x) \ne x$ and $t(z) \ne z$
for $x,z \in X$.
Since $x \in t^{-1}(t(x))$ and $z \in t^{-1}(t(z))$,
we find $t(x) = t(z) = y$.

$\ref{prelim:prop:item:mapping} \Rightarrow~\ref{prelim:prop:item:invimg2}$: 
We can assume $t(x) \ne x$ for some $x \in X$, as otherwise~$t$
equals the identity transformation and the claim obviously holds.
Then set $y = t(x)$. Next, consider $z \in X$
with $z \notin t^{-1}(y)$.
By assumption, $t(z) = z$. Now, assume $t(z) = t(z')$ for some $z' \in X$
with $z \ne z'$.
Then $t(z') \ne z'$ yields $t(z') = y$. So, this is not possible
and we have $t^{-1}(z) = \{z\}$.

$\ref{prelim:prop:item:invimg1a} \Rightarrow~\ref{prelim:prop:item:invimg2}$: 	Let $z \in X$ with $z \ne y$
and assume $z \notin t^{-1}(y)$.
Now, for $x \in X$ with $x \in t^{-1}(z)$, by assumption, we get $z = t(x) = x$
and so $t^{-1}(z) \subseteq \{z\}$.

$\ref{prelim:prop:item:invimg2} \Rightarrow~\ref{prelim:prop:item:excluded2}$: 
 Let $z \in t(X)$ and suppose $z \ne y$.
 We have $z = t(x)$ for some $x \in X$.
 Thus $x \in t^{-1}(z)$. Since $z \ne y$,
 by assumption $t^{-1}(z) \subseteq \{z\}$
 and so $z = x$. This gives $t^{-1}(z) = \{z\}$.
 Since $t^{-1}(z) = \{z\}$ for $z \in t(X)$ with $z \ne y$,
 and the union of the inverse images equals~$X$,
 we have that the remaining points $\excl{t} \cup \{y\}$
 must be in the inverse image of $t^{-1}(y)$, 
 and since this union is pairwise disjoint, no other point is in $t^{-1}(y)$  
 and so $t^{-1}(y) = \excl{t} \cup \{y\}$.
 
$\ref{prelim:prop:item:excluded2} \Rightarrow \ref{prelim:prop:item:inverse_super_or_subset2}$: Let $A \subseteq X$.
Then
\[
t^{-1}(A) = t^{-1}(A \cap t(X))
= \begin{cases}
	A \cap t(X) & \text{if } y \notin A, \\
	\excl{t} \cup A & \text{if } y \in A. 	
\end{cases}
\]
So, $t^{-1}(A) \subseteq A$ for $y \notin A$
and $A \subseteq t^{-1}(A)$ for $y \in A$.
	
$\ref{prelim:prop:item:inverse_super_or_subset2} \Rightarrow~\ref{prelim:prop:item:one-attractor}$: 
 For $z \in X$, consider $A = \{t(z)\}$. Then $A \subseteq t^{-1}(A)$ or $t^{-1}(A) \subseteq A$ and as $t^{-1}(A) \ne \emptyset$,
 we get $t(z) \in t^{-1}(t(z))$ for both cases.
 Hence, by~\crefitem{prelim:prop:idempotent}{prop:idemp:item:inv},~$t$ is idempotent. Then, for $\dupl{t} = \emptyset$ we have $t = \operatorname{id}_X$.
 Next, for $\dupl{t} \ne \emptyset$,
 by \crefitem{prelim:cor:idempotent}{prelim:idemp:5}
 there exists $y \in \dupl{t}$
 such that  $x \in t^{-1}(y) \cap \excl{t}$
 for some $x \in X$.
 Let $y' \in \dupl{t}$ and set $A = \{x,y'\}$.
 Then $t^{-1}(A) = t^{-1}(y')$ and, as $t^{-1}(y')$ contains at least two points,
 for $A \subseteq t^{-1}(y')$ or $t^{-1}(y') \subseteq A$
 we get $x \in t^{-1}(y')$ in both cases.
 Now $x \in t^{-1}(y) \cap t^{-1}(y')$ gives $y = y'$
 and so~$t$ is a one-point singular.
%
% Alternativer Beweis für one-point singular:
%	Suppose $y,y' \in \dupl{t}$.
%	Then there exist $z,z' \in X$ with $z \in t^{-1}(y)$, 
%	$z' \in t^{-1}(y')$, $z \ne y$ and $z' \ne y'$. % z,z' excluded
%	Set $A = \{y, z'\}$. 
%	First, suppose $A \subseteq t^{-1}(A)$.
%	Then $z' \in t^{-1}(A)$.
%	As the inverse images partition $X$,
%	we must have $y' \in A$.
%	As $y' \ne z'$, this implies $y = y'$.
%	Next, suppose $t^{-1}(A) \subseteq A$.
%	As $y \in A$ we have $z \in t^{-1}(A)$.
%	So, $z \in A$.
%	Now, $z \ne y$ implies $z = z'$,
%	which yields $y = y'$.
%   Hence $t$ is a one-point singular.
\end{proof}

Regarding \crefitem{prelim:thm::one-attractor-idempotent}{prelim:prop:item:equation},
inductively, from \cref{prelim:trans:eqn:SI-one-singular0} (but it is also easy to verify directly)
we get the following identity:
\begin{equation} \label{prelim:eqn:concat-SI}
\concat{\SI{a_1}{a_2}}{\concat{\cdots \SI{a_2}{a_3}}{\SI{a_{m-1}}{a_m}}}
= [a_{m-1}\mapsto a_m] \cdots [a_2 \mapsto a_m][a_1 \mapsto a_m]. 
\end{equation}

Next, we state a new characterization of idempotents on finite sets whose duplicate set contains at most one point.

\begin{proposition} 
\label{prelim:prop:one-attractor-idempotent-fin}
  Let $t \colon [n] \to [n]$. 
  Then $t = \operatorname{id}_{[n]}$ or $t$ is an idempotent one-point singular
  if and only if 
  for all $A \subseteq [n]$
  we have that $|\applyfunc{A}{t^{-1}}| = |A|$
  implies $\applyfunc{A}{t^{-1}} = A$.		
\end{proposition}
\begin{proof}
	We can assume $t \ne \operatorname{id}_{[n]}$.
	First, suppose $t$ is an idempotent one-point singular. Let $A \subseteq [n]$ be such that
	$|t^{-1}(A)| = |A|$.
	By~\crefitem{prelim:thm::one-attractor-idempotent}{prelim:prop:item:inverse_super_or_subset2} we have $t^{-1}(A) \subseteq A$ or $A \subseteq t^{-1}(A)$.
	Finiteness implies $t^{-1}(A) = A$ in both cases.
		
	Conversely, consider a set $A \subseteq [n]$.
	Since $\defect{t} + |A \cap t([n])| \ge |t^{-1}(A \cap t([n]))|$ by \crefitem{prelim:prop:gentransdef}{prelim:prop:item:defA},
	we can choose $B \subseteq \excl{t}$
	such that $|B| + |A \cap t([n])| = |t^{-1}(A \cap t([n]))|$
	where we let $B \subseteq A \cap \excl{t}$ if $|B| \le |A \cap \excl{t}|$
	and $A \cap \excl{t} \subseteq B$ otherwise.
	Since $B \subseteq \excl{t}$
	we have $|B| + |A \cap t([n])| = |B \cup (A \cap t([n]))|$
	and $t^{-1}(A \cap t([n])) = t^{-1}(( A \cap t([n]) ) \cup B)$.
	Thus $| t^{-1}(( A \cap t([n]) ) \cup B) | = | ( A \cap t([n]) ) \cup B) |$
	and so $t^{-1}( ( A \cap t([n]) ) \cup B) = ( A \cap t([n]) ) \cup B$.
	Since $t^{-1}(A) = t^{-1}(A \cap t([n]))$
	and $A = (A \cap t([n])) \cup (A \cap \excl{t})$,
	we find $t^{-1}(A) \subseteq A$ for $B \subseteq A \cap \excl{t}$
	and $A \subseteq t^{-1}(A)$
	for $A \cap \excl{t}\subseteq B$.
	Therefore, by \crefitem{prelim:thm::one-attractor-idempotent}{prelim:prop:item:inverse_super_or_subset2}, 
	it follows that $t$ is an idempotent one-point singular.
\end{proof}

\begin{proposition} \label{prelim:trans:idempt-one-point:prop:comp} 
	Let $t, h \colon X \to X$ be idempotent one-point singulars. Then the following are equivalent.
	\begin{enumerate}
		\item \label{prelim:trans:idempt-one-point:prop:comp:1} $\concat{t}{h}$ is an idempotent one-point singular.
			
		\item \label{prelim:trans:idempt-one-point:prop:comp:2a} 
		$\move{t} \cap h(\move{h} \cap t(X)) \subseteq \fix{\concat{t}{h}}$
		and 
		$\dupl{\concat{t}{h}} = h(\dupl{t})$.\footnote{Recall that
			$\dupl{h_{|t(X)}} \subseteq h(\dupl{t})$ is equivalent to $\dupl{\concat{t}{h}} = h(\dupl{t})$
			by \cref{prelim:eqn:dupl0}.}
		
		\item \label{prelim:trans:idempt-one-point:prop:comp:2} $h(\move{h} \cap t(X)) \subseteq h(\dupl{t})$.
				
		\item \label{prelim:trans:idempt-one-point:prop:comp:3} $t(X) \subseteq \fix{h}$ or $\dupl{h} = h(\dupl{t})$.\footnote{This is equivalent to the statement that $\move{h} \cap t(X) \ne \emptyset$ implies $\dupl{h} = h(\dupl{t})$.}
	    
	    \item \label{prelim:trans:idempt-one-point:prop:comp:4} We have $\concat{t}{h} = t$
	    or with $\move{t} \cup \move{h} = \{a_1, \ldots, a_{m-1}\}$ and $\dupl{h} = \{a_m\}$
	    for distinct $a_1, \ldots, a_m \in X$
	    we have
	    \[
	    \concat{t}{h} = [a_1 \mapsto a_m] \cdots [a_{m-1} \mapsto a_m].
	    \]
\end{enumerate}	
\end{proposition}
\begin{proof}
$\ref{prelim:trans:idempt-one-point:prop:comp:1} \Rightarrow \ref{prelim:trans:idempt-one-point:prop:comp:2a}$:
 The first inclusion follows from \cref{prelim:cor:comp-idempotents-idempotent}.
 Since $\concat{t}{h}$ and $t$ are one-point singulars, we must have $\dupl{\concat{t}{h}} = h(\dupl{t})$.
 	
$\ref{prelim:trans:idempt-one-point:prop:comp:2a} \Rightarrow \ref{prelim:trans:idempt-one-point:prop:comp:2}$: 
Let $x \in X$ be such that $h(t(x)) \ne t(x)$.
If $t(h(t(x))) \ne h(t(x))$, then, by assumption, $h(t(h(t(x)))) = h(t(x))$
and thus $h(t(x)) \in \dupl{h_{|t(X)}}$.
If $t(h(t(x))) = h(t(x))$, then $h(t(x)) \in t(X)$.
By idempotency of~$h$, we find $h(h(t(x))) = h(t(x))$.
Thus $h(t(x)) \in \dupl{h_{|t(X)}}$.
Since $\dupl{h_{|t(X)}} \subseteq \dupl{\concat{t}{h}} = h(\dupl{t})$
we obtain the claim.
 
$\ref{prelim:trans:idempt-one-point:prop:comp:2} \Rightarrow \ref{prelim:trans:idempt-one-point:prop:comp:3}$: 
Let $t(x) \in \move{h}$.
Since~$h$ is an idempotent, we have $h(\move{h}) = \dupl{t}$. Since it is a one-point singular, we find 
$\dupl{h} = \{ h(t(x)) \}$.
By assumption, we have $h(t(x)) \in h(\dupl{t})$. Since~$t$ is a one-point singular
we get $h(\dupl{t}) = \{ h(t(x)) \}$.
Thus $\dupl{h} = h(\dupl{t})$.

%$\ref{prelim:trans:idempt-one-point:prop:comp:2} \Rightarrow \ref{prelim:trans:idempt-one-point:prop:comp:2a}$:
%It suffices to show $\dupl{h_{|t(X)}} \subseteq h(\dupl{t})$.
%Let $z \in \dupl{h_{|t(X)}}$.
%Since two points from~$t(X)$ are mapped to~$z$,
%we have $z = h(t(x))$ for some $x \in X$ with $t(x) \ne z$.
%By assumption, this gives $z \in h(\dupl{t})$.

$\ref{prelim:trans:idempt-one-point:prop:comp:3} \Rightarrow \ref{prelim:trans:idempt-one-point:prop:comp:4}$: 
Suppose $ht \ne t$, i.e., $t(X) \nsubseteq \fix{h}$.
Then, by \cref{prelim:trans:idempt-one-point:prop:comp:3}, we have $\dupl{t} \subseteq h^{-1}(a_m)$.
Let $x \in \move{t}$. Then $t(x) \in \dupl{t}$ since $t$ is idempotent
and so $h(t(x)) = a_m$.
Next, let $x \in \fix{t}$ and $x \in \move{h}$.
Then $h(t(x)) = h(x) = a_m$ since~$h$ is idempotent.
Next, for $x \in \fix{t}\cap \fix{h}$
we have $h(t(x)) = x$.
Thus, precisely the points from $\move{t} \cup \move{h}$ are mapped to~$a_m$
and all other points are fixed. This is precisely
the transformation described by $[a_1 \mapsto a_m] \cdots [a_{m-1} \mapsto a_m]$.\footnote{Alternatively, the equation for $\concat{t}{h}$ follows by deriving $\dupl{\concat{t}{h}} = \dupl{h}$,
	substituting this into \cref{prelim:trans:idemp:prop:inv-dupl},
	and noting that all points outside of $t^{-1}(\dupl{t}) \cup h^{-1}(\dupl{h}) = \move{t} \cup \move{h} \cup \{a_m\}$
	are fixed by $\concat{t}{h}$.}

$\ref{prelim:trans:idempt-one-point:prop:comp:4} \Rightarrow 
\ref{prelim:trans:idempt-one-point:prop:comp:1}$: Clearly, in each case,
$\concat{t}{h}$ is an idempotent one-point singular.
\end{proof}

\begin{corollary}
	Let $t, h \colon X \to X$ be idempotent one-point singulars
	such that $\concat{t}{h}$ is an idempotent one-point singular.
	Then
	\[
	(\concat{t}{h})^{-1}(\dupl{\concat{t}{h}})
	=
	\begin{cases}
		t^{-1}(\dupl{t}) & \text{if } t(X) \subseteq \fix{h}, \\
		t^{-1}( \dupl{t} ) \cup h^{-1}( \dupl{h} ) & \text{otherwise,}				
	\end{cases}
	\]
	where we also have $t^{-1}( \dupl{t} ) \cup h^{-1}( \dupl{h} )
	= t^{-1}( \dupl{t} ) \cup ( t(X) \cap h^{-1}( \dupl{h} ))$.
\end{corollary}
\begin{proof}
	First, suppose $t(X) \subseteq \fix{h}$.
	Then $\concat{t}{h} = t$ and $\dupl{\concat{t}{h}} = \dupl{t}$,
	which gives the equality for this case.
	
	Next, suppose $t(X) \nsubseteq \fix{h}$. By \cref{prelim:trans:idempt-one-point:prop:comp},
	we have $\dupl{h} = \dupl{\concat{t}{h}}$.
	Then the claim follows from \cref{prelim:trans:idemp:prop:inv-dupl}.
\end{proof}

\begin{remark}
	\Cref{prelim:trans:fig:duplconcatingduplf} 
	shows two idempotent one-point singulars
	whose composition results in an idempotent one-point singular.
	The first two pictures on the left show the case $t(X) \subseteq \fix{h}$ (see also \cref{prelim:trans:idemp:prop:tX-in-fixh}),
	the rightmost picture shows the case $\dupl{h} = h(\dupl{t})$.
\end{remark} 

\begin{corollary} \label{prelim:cor:concat_SI} 
	Let $t, h \colon X \to X$ be two simple idempotents.
	If $\concat{t}{h}$ is a simple idempotent,
	then $\concat{t}{h} \in \{t,h\}$.	
\end{corollary}
\begin{proof}
	By \crefitem{prelim:trans:idempt-one-point:prop:comp}{prelim:trans:idempt-one-point:prop:comp:4},
	if $\concat{t}{h}$ is a simple idempotent,
	we must have $\concat{t}{h} = t$
	or $\move{t} \cup \move{h} = \{a\}$ and $\dupl{h} = \{b\}$,
	which gives $h = \SI{a}{b}$,
	and $\concat{t}{h} = \SI{a}{b}$.
\end{proof}

\begin{remark}
	Several simple idempotents composed
	can result in new simple idempotents,
	a result that fails for two simple idempotents
	by \cref{prelim:cor:concat_SI}.
	For example, consider $t_1 = \SI{1}{2}$, 
	$t_2 = \SI{2}{3}$, and $t_3 = \SI{3}{1}$.
	Then $\concat{\concat{t_3}{\concat{t_2}{t_1}}}{\concat{t_3}{\concat{t_2}{t_1}}} = \SI{1}{3}$
	where $\concat{\concat{t_2}{t_1}}{\concat{t_3}{\concat{t_2}{t_1}}} = [1 \mapsto 3 \mapsto 2]$.
\end{remark}

\begin{remark}
Let $t,h \colon X \to X$ such that $t$ (resp.~$h$) is a simple idempotent
moving $a$ to the distinct point~$b$.
Then a simple case yields that $\concat{t}{h}$ is a simple idempotent
if and only if $\move{h} \subseteq \{a,b\}$ and $h(b) \in \{a,b\}$ 
(resp.~$\move{t} \subseteq \{a,b\}$ and $t(b) \in \{a,b\}$).
In this case we have that $\concat{t}{h}$ equals $\SI{a}{b}$ or $\SI{b}{a}$ (resp.~$\SI{a}{b}$ for $t(a) \in \{a,b\}$ 
and $\SI{a}{t(a)}$ for $t(a) \notin \{a,b\}$). The next example illustrates
some cases.
\end{remark}

\begin{example}
	Let $t, h \colon X \to X$ for $X = \{1,2,3\}$
	with $t(1) = t(2) = 2$, $t(3) = 1$, 
	and	$h(1) = h(3) = 3, h(2) = 1$. 
	Then we have
	\begin{enumerate}
		\item $\concat{t}{\SI{1}{3}} = \SI{1}{2}$.		
		
		\item $\concat{\SI{1}{2}}{h} = \SI{2}{1}$.
		
		\item $\concat{\SI{1}{2}}{\SI{1}{3}} = \SI{1}{2}$.	
	\end{enumerate}
\end{example}